\section*{Scriptural Date and Location}

{\footnotesize Seven passages are directly appointed; the three orations quote
no Scripture and appear nowhere below. Psalms carry the missal's Vulgate
number, the modern number in parentheses, with the convention named where the
two English series part. \textbf{A traditional attribution's regnal years are
the era of the author a psalm title names, and are neither a date of
composition nor an occasion}: the corpus keeps attribution, setting and
composition apart, and the Communion cell prints all three at once.\par}

\begin{dossiertable}
Offertory & Ps. 39:14--15 (Hebr. 40:14--15; Eng. 40:13--14) & Title names David, no place or occasion & \chronodate{offertory}{\chronologyannotation{offertory}}\\
\cmidrule(lr){1-4}
\dossierprose{The Clementine title verse reads \latin{In finem. Psalmus ipsi
David}, the Douay \emph{Unto the end, a psalm for David himself}; \textbf{no
place, episode or audience is named there or anywhere in the psalm}, and beside
the attributed reign the corpus adds only the boundary shared by the whole
Psalter, whose own source declines to date any individual psalm securely.
\textbf{Ps.~39 is one of the two psalms here whose numbered Hebrew title an
English Bible leaves unnumbered}, so an English psalter at 40:14--15 gives the
Offertory's second verse and the one after it.}
\midrule
Communion & Ps. 70:16-17 et 18 (Eng. 71:16--18) & Title: the sons of Jonadab, the former captives & \chronodate{communion}{\chronologyannotation{communion}}\\
\cmidrule(lr){1-4}
\dossierprose{\textbf{Three settings stand here at once, below the attributed
reign, and the corpus resolves none of them.} The Clementine title reads
\latin{Psalmus David, filiorum Jonadab, et priorum captivorum}; the flight from
Absalom reaches this psalm through the corpus's grouping of Ps.~62 with Ps.~70
and is asserted at Ps.~3's superscription, not at this psalm's; the third year
of Joakim comes from Daniel~1:1; and the year of the world beside it is the
Ussher-attributed marginal apparatus of the Haydock printing, carried as a
reported figure and excluded from the corpus's own basis. \textbf{The
composition boundary here is from a critical-profile source; every other claim
on this formulary is from the traditional profile.}}
\midrule
Introit & Ps. 85:3, 5, with v.~1 (Eng. 86:3, 5, 1) & Title names David, no occasion & \chronodate{introit}{\chronologyannotation{introit}}\\
\cmidrule(lr){1-4}
\dossierprose{The Clementine heads the psalm \latin{Oratio ipsi David} inside
v.~1 --- the words the Introit's psalm verse begins after --- and \textbf{names
no place, no episode and no audience}. Beside the attributed reign the corpus
carries composition only for these verses, at the whole-Psalter boundary above.
Augustine's single sermon on this psalm asks no chronological question of it,
and his lemma at v.~1 reads \latin{egenus et inops ego sum} where the missal
sings \latin{inops et pauper sum ego}.}
\midrule
Alleluia & Ps. 97:1 (Eng. 98:1) & Title names David, no place & \chronodate{alleluia}{\chronologyannotation{alleluia}}\\
\cmidrule(lr){1-4}
\dossierprose{The Clementine folds \latin{Psalmus ipsi David} into v.~1 and the
chant drops it, opening at \latin{Cantate Domino}. \textbf{The two Nativity
labels in the Date cell are prophetic-referent claims and not composition
claims, and neither is made of Ps.~97 alone}: both reach this verse through the
corpus's grouping of Pss.~67, 95, 96 and 97, and both stand disputed side by
side. The composition side carries only the whole-Psalter boundary.}
\midrule
Gradual & Ps. 101:16--17 (Hebr. 102:16--17; Eng. 102:15--16) & Title: a poor man; no name, no place & \chronodate{gradual}{\chronologyannotation{gradual}}\\
\cmidrule(lr){1-4}
\dossierprose{The title verse reads \latin{Oratio pauperis, cum anxius fuerit,
et in conspectu Domini effuderit precem suam} --- a genre and a condition, with
\textbf{no name, no place and no occasion}. \textbf{Ps.~101 is the one appointed
psalm the attribution to David does not reach}, and the corpus carries
composition only. \textbf{This is also the second psalm whose numbered Hebrew
title an English Bible leaves unnumbered}, so an English psalter at 102:16--17
gives the Gradual's second verse and the one after it. Augustine expounds both
verses in \work{Enarratio in Ps.~101} sermo~I §§16--17 on a lemma reading a
future \latin{aedificabit} and \latin{in gloria sua}.}
\midrule
Gospel & Luke 14:1--11 & Written by Luke; no place, no audience & \chronodate{gospel}{\chronologyannotation{gospel}}\\
\cmidrule(lr){1-4}
\dossierevent{Event: a sabbath meal in the house of one of the chief of the
Pharisees (14:1), the third and last of the sabbath-healing controversies in
Luke after 6:6--11 and 13:10--17, inside the journey narrative and closing
immediately before the parable of the great supper (14:16--24). \textbf{The
corpus makes no dated assertion about this sabbath dinner}: it carries no
narrated-event claim for the pericope at all.}
\cmidrule(lr){1-4}
\dossierprose{\textbf{The two composition labels stand disputed side by side}:
one from the \work{Catholic Encyclopedia}, vol.~XIV, New Advent article
14530a, the other from the Pontifical Biblical Commission's responsa \work{De
auctore, tempore et veritate evangeliorum Marci et Lucae}, which fixes the
writing relative to Acts and to the Roman imprisonment rather than by a year.
\textbf{Neither source names a place of writing, a first audience or a life
stage}, and they are the nearest that would carry them for this Gospel.
\textbf{The corpus returns no critical-profile claim for this Gospel.}}
\midrule
Epistle & Eph. 3:13--21 & Headed to the Ephesians; no place named & \chronodate{epistle}{\chronologyannotation{epistle}}\\
\cmidrule(lr){1-4}
\dossierprose{The missal heads the lesson \latin{Léctio Epístolæ beáti Pauli
Apóstoli ad Ephésios}, and the pericope names its writer's imprisonment only at
\latin{tribulatiónibus meis pro vobis}. \textbf{Both labels in the Date cell are
disputed and both are traditional-profile}, from the \work{Catholic
Encyclopedia}, vol.~V, New Advent article 05485a, and vol.~XI, article 11567b;
\textbf{the corpus returns no critical-profile claim for this letter}.
\textbf{Neither article names a place of writing or a first audience}, and they
are the nearest sources that would carry them.}
\end{dossiertable}

\clearpage
